\documentclass[../../main.tex]{subfiles}
\begin{document}

{\bf[解]} $x'=2t$，$y'=2t+1$であるから，下表を得る．

\begin{table}[htb]
\centering
\caption{$(x,y)$ の増減表}
\begin{tabular}{c|ccccc}
$t$ & $\cdots$ & $-\dfrac{1}{2}$ & $\cdots$ & $0$ & $\cdots$ \\
\hline
$x'$ & $-$ & $-$ & $-$ & $0$ & $+$ \\
\hline
$y'$ & $-$ & $0$ & $+$ & $+$ & $+$ \\
\hline
$(x,y)$ & $\swarrow$ & $\left(\dfrac{5}{4},\,-\dfrac{9}{4}\right)$ & $\nwarrow$ & $(1,\,-2)$ & $\nearrow$
\end{tabular}
\end{table}

また，極限値は，$t\to\pm\infty$の時，$x,y\to\infty$である．ゆえにグラフの概形は下図のようになる．
% \begin{center}
%      \scalebox{1}{\input{ut-61-5p}}
% \end{center}

\begin{figure}[htb]
\centering
\begin{tikzpicture}[scale=1.1, >=stealth, every node/.style={font=\small}]

  % --- 1. 積分領域（面積 S）の塗りつぶし ---
  % t \in [-2, 1] の区間。
  % t = -2: (5, 0)
  % t = 1 : (2, 0)
  % 下側曲線 y_- (t: -2 -> 0) と上側曲線 y_+ (t: 0 -> 1)、および x軸で囲まれた閉領域
  \fill[pattern=north east lines, pattern color=black!30]
    plot[domain=-2:1, samples=60] ({(\x)^2 + 1}, {(\x)^2 + \x - 2})
    -- (2, 0) -- (5, 0) -- cycle;

  % --- 2. 座標軸 ---
  \draw[->, thick] (-0.5, 0) -- (6.0, 0) node[right] {$x$};
  \draw[->, thick] (0, -3.0) -- (0, 2.5) node[above] {$y$};
  \node[below left] at (0,0) {$O$};

  % --- 3. 媒介変数曲線 (x(t), y(t)) の描画 ---
  % x(t) = t^2 + 1, y(t) = t^2 + t - 2
  \draw[very thick, blue!80!black, domain=-2.3:1.4, samples=100]
    plot ({(\x)^2 + 1}, {(\x)^2 + \x - 2});

  % --- 4. 進行方向を示す矢印 (t の増加方向) ---
  % t = -1 付近: (2, -2) 方向へ
  \draw[->, thick, blue!80!black] ({(-1.1)^2 + 1}, {(-1.1)^2 - 1.1 - 2}) 
    -- ({(-0.9)^2 + 1}, {(-0.9)^2 - 0.9 - 2});
  % t = 0.5 付近: (1.25, -1.25) 方向へ
  \draw[->, thick, blue!80!black] ({(0.4)^2 + 1}, {(0.4)^2 + 0.4 - 2}) 
    -- ({(0.6)^2 + 1}, {(0.6)^2 + 0.6 - 2});

  % --- 5. 特徴的な点のプロットと補助点線 ---
  % (i) t = -2: (5, 0) [x軸との交点]
  \fill (5, 0) circle (1.3pt) node[above] {$(5,0)$};
  \node[above right, font=\footnotesize, gray!80!black] at (5, 0) {($t=-2$)};

  % (ii) t = -1/2: (5/4, -9/4) [y の極小点]
  \coordinate (Pmin) at (1.25, -2.25);
  \fill[red!80!black] (Pmin) circle (1.3pt);
  \draw[dashed, gray!70] (1.25, 0) -- (Pmin) -- (0, -2.25);
  \node[below=2pt, font=\footnotesize] at (1.25, 0) {$\frac{5}{4}$};
  \node[left=2pt, font=\footnotesize] at (0, -2.25) {$-\frac{9}{4}$};
  \node[below right, red!80!black, font=\footnotesize] at (Pmin) {$\left(\frac{5}{4},\,-\frac{9}{4}\right)$};

  % (iii) t = 0: (1, -2) [x の最左端 / 停留点]
  \coordinate (Pleft) at (1, -2);
  \fill[red!80!black] (Pleft) circle (1.3pt);
  \draw[dashed, gray!70] (1, 0) -- (Pleft) -- (0, -2);
  \node[below=2pt, font=\footnotesize] at (1, 0) {$1$};
  \node[left=2pt, font=\footnotesize] at (0, -2) {$-2$};
  \node[left=2pt, red!80!black, font=\footnotesize] at (Pleft) {$(1,-2)$};

  % (iv) t = 1: (2, 0) [x軸との交点]
  \fill (2, 0) circle (1.3pt) node[above left] {$(2,0)$};
  \node[above, font=\footnotesize, gray!80!black] at (2.1, 0.15) {($t=1$)};

  % --- 6. ラベル ---
  % 曲線部分の識別
  \node[blue!80!black, above left] at (3.5, 0.8) {$y_+\ (t \ge 0)$};
  \node[blue!80!black, below] at (3.5, -1.2) {$y_-\ (t \le 0)$};

  % 求める面積 S の領域ラベル
  \node[fill=white, inner sep=2pt, font=\bfseries] at (2.6, -0.9) {$S$};

\end{tikzpicture}
\caption{媒介変数曲線 $x = t^2 + 1,\ y = t^2 + t - 2$ の概形と求める領域 $S$}
\end{figure}

ここでグラフの下側を$y_-$，上側を$y_+$とすると，求める面積$S$は，
\begin{align}
     S&=\int_1^2(y_+-y_-)dx-\int_2^5y_-dx \\
     &=\int_0^1y_+\dfrac{dx}{dt}dt-\int_0^{-1}y_-\dfrac{dx}{dt}dt-\int_{-1}^{-2}y_-\dfrac{dx}{dt}dt \\
     &=\int_{-2}^1y\dfrac{dx}{dt}dt \\
     &=\int_{-2}^1(t^2+t-2)2tdt \\
     &=2\left[\dfrac{1}{4}t^4+\dfrac{1}{3}t^3-t^2\right]_{-2}^1=\dfrac{9}{2}
\end{align}
である．$\cdots$(答)

{\bf[別解]}

$xy'-x'y=-t^2+6t+1$より，ガウスグリーンの定理から，
\begin{align}
     T=\dfrac{1}{2}\int_{-2}^1-(-t^2+6t+1)dt=\dfrac{9}{2}
\end{align}     
である．$\cdots$(答)     
\newpage
\end{document}